The analysis of gunshot wounds (Entry and Exit) and the estimation of the firing distance (MLSD - Medical-Legal Shooting Distance) are challenging forensic tasks. They typically demand highly specialized experts and depend on sub-optimal image capture conditions at crime scenes. The advent of Deep Learning has brought new opportunities to automate and assist in these tasks, yet it still faces difficulties with limited and imbalanced forensic datasets.
This work investigates the use of a hybrid architecture called TransferKAN for the classification of ballistic wounds on real crime scene images from the FDCPUnBGunshotDB dataset. The proposed architecture employs a pre-trained convolutional neural network (ResNet152) to extract robust spatial features, replacing the traditional dense classifier with a Kolmogorov-Arnold Network (KAN), which is capable of modeling complex non-linear relationships more efficiently.

The approach was evaluated under a rigorous k-fold cross-validation scheme to prevent data leakage, coupled with advanced Data Augmentation techniques to mitigate the inherent class imbalance of forensic samples. The performance of the TransferKAN model was compared with results from literature (a 2024 paper baseline), demonstrating technical superiority.

For the wound type classification task (Entry vs. Exit), the TransferKAN model with k-fold validation achieved an Accuracy of 86.0\% and a weighted AUC of 0.908, with a weighted F1-Score of 0.862. In the firing distance classification (MLSD), the TransferKAN showed improvements in weighted precision, weighted recall, and balanced F1-score, reaching 87.3\% Accuracy. These findings demonstrate that our TransferKAN approach with k-fold validation achieved the best performance and overall results, highlighting the potential of using KAN networks as classifiers to assist forensic pathology in real-world scenarios.
